\section{案例分析与下一步}

\subsection{Arc AGI Puzzle Walkthrough}
在 Arc AGI 这条 benchmark 曲线上，我们可以用以下步骤驱动 inference loop：1) Prompt 里提前声明这个 puzzle 是 high-difficulty；2) Sampling 阶段生成 30~40 个 candidate response；3) Search 阶段把这些 candidate 当作 tree seed，用 stepwise scorer 计算 promisingness 并且 intro probability threshold；4) Iteration 阶段运行 trace feedback，自我检查每条 response 的 reasoning chain，并借助 verifier margin 提示多数票选择。整个流程也会记录 every candidate’s token cost 与 latency，供 future tuning。

\begin{importantbox}{Arc AGI Pipeline 观察}
把 Arc AGI 视为全流程演练场：prompt → sampling → search → trace → voting，每一环都需要 instrumentation；如果 token cost 突然翻倍，看 sampling diversity 与 verifier margin 是否崩溃，然后用 runbook 回滚到 conservative mode。
\end{importantbox}

\subsection{高价值任务的推理时部署}
对于像 AlphaProof 这种竞争性 Code/Math task，pipeline 里通常还会加入 extra verification: 1) Run deterministic verifier (e.g. proof checker) as soon as the candidate completes; 2) If verification fails, reroute to alternative candidate with similar margin; 3) Log the mismatch and feed trace back to sampling controller for prompt iteration。这样的流程使 inference-time compute 成为一条 closed-loop control rather than a one-shot guess。

\begin{knowledgebox}{高价值部署要点}
1. 把 deterministic verifier 作为 fallback，避免 majority vote 把 wrong candidate 当成正确答案；2. 用 trace replay 复现每一次 verification failure，以便 prompt 出新的 examples；3. 把 close-loop pipeline 的 metrics（token cost, verifier success, trace iteration）记录在 performance dashboard。
\end{knowledgebox}

\subsection{本章小结}
案例分析展示了如何把推理时 compute pipeline 实际落地到 Arc AGI 与 AlphaProof 类任务：先确定 prompt 级别，再让 sampling/search/trace 组成控制钮，最后把 verification 反馈写入 iteration loop 以持续改进。

